\section{The model layer: part tree, placement, and mass properties}\label{sec:model}

The model layer describes a rocket as a composite tree of \code{model::part::Part} nodes, each of
which is simultaneously one component and the root of a sub-tree whose mass, center of mass, and
inertia tensor it aggregates on demand (\src{model/parts/Part.h}{25}). Three mechanisms define the
layer: an ownership tree with move-only attachment, stable ids, and polymorphic deep-copy; an
intent-based placement model whose only stored spatial relationship is a \code{StationLink}, every
absolute pose being derived by a single resolver; and a two-pass, time-aware mass-properties walk
cached behind a mass-delta gate. The layer also exhibits, in concentrated form, this review's
recurring motif: several of its seams --- the reserved placement quaternion, the per-seam radial
check, the live-length re-resolution machinery, the Barrowman composition --- are finished and
tested but have no production consumer yet.

\subsection{Tree ownership, identity, and duplication}\label{sec:model-tree}

\code{Part} is an abstract base --- \code{typeName()} and \code{cloneShallow()} are pure
(\src{model/parts/Part.h}{181}, \src{model/parts/Part.h}{217}) --- owning its children as a vector
of \cppinline|std::pair<std::shared_ptr<Part>, StationLink>| (\src{model/parts/Part.h}{288}),
with a non-owning \cppinline|Part* parent| used solely to propagate staleness upward
(\src{model/parts/Part.h}{222}). At the value level a part is neither copyable nor movable ---
assignment is deleted and the copy constructor protected, so a subclass cannot be sliced and ``a
part lives at one place in one tree'' (\src{model/parts/Part.h}{66}). Attachment transfers
ownership by move: \code{addChildPart} rejects a null child, an already-parented child, and any
cycle-forming attachment as logged no-ops (\src{model/parts/Part.cpp}{73}), then re-parents,
emplaces, and dirties both the composite and placement caches up the ancestor chain
(\src{model/parts/Part.cpp}{97}). \code{removeChildById} moves the owning pointer out, clears the
detached node's parent, dirties the ex-parent chain, and returns the intact sub-tree
(\src{model/parts/Part.cpp}{277}).

Identity is separated from naming: each part gets a process-unique \code{Id} from a relaxed atomic
counter starting at 1, with 0 reserved as none/invalid (\src{model/parts/Part.cpp}{28}); ids
survive adoption, are refreshed on copy, and are deliberately not serialized
(\src{model/parts/Part.h}{171}). Duplication has one path: \code{clone()} calls the pure
\code{cloneShallow()} (implemented per type via the protected copy constructor,
\src{model/parts/BodyTube.h}{63}) to reproduce the node's dynamic type, then recursively clones
children and copies each \code{StationLink} verbatim (\src{model/parts/Part.cpp}{106}). One
consequence deserves stating plainly: apart from \code{setMass}/\code{setI}
(\src{model/parts/Part.h}{74}), a part is immutable after construction --- no concrete type
exposes a geometry setter (\src{model/parts/BodyTube.h}{69}), so changing a dimension means
building a replacement part and losing the id.

\subsection{Frame convention and the StationLink placement model}\label{sec:model-placement}

Every part lives in its own \zfwd{} frame, occupying $z \in [-L, 0]$ with the fore plane at the
local origin (\src{model/parts/Part.h}{130}). Geometry is exposed through closed-form silhouette
virtuals \code{radiusOuterAt}/\code{radiusInnerAt} (\src{model/parts/Part.h}{138}), an
\code{axialLength()} envelope span defaulting to \code{getLength()}
(\src{model/parts/Part.h}{146}), an \code{isSolid()} heuristic (\src{model/parts/Part.h}{151}),
and \code{stationAt(station01)}, which clamps to $[0,1]$ and maps the fraction to $z = (s-1)L$
(\src{model/parts/Part.cpp}{305}). \code{innerCapacityAt} encodes the solid-host rule the
diagnostics rely on: a bored part is bounded by its inner wall, a solid part by its outer skin
(\src{model/parts/Part.cpp}{314}).

Placement is expressed as intent, not coordinates. A \code{StationLink} carries a fractional
station on the parent (default 0, aft plane), one on the child (default 1, fore plane), a gap in
metres whose sign convention is selected by a \code{SeatKind}, and a quaternion \code{childRot}
``reserved for 6-DOF; identity in 3-DOF'' (\src{model/parts/Placement.h}{49}). \code{SeatKind} is
closed at three values --- \code{Abut} (rim-to-rim, gap a forward standoff), \code{NestInBore}
(child OD inside parent ID, gap an insertion depth moving the child aft), and \code{OnSurface}
(child on the parent's outer wall) --- since a rigid axisymmetric stack admits only three radial
relationships (\src{model/parts/Placement.h}{32}). Sugar verbs \code{abut}, \code{nestInBore}, and
\code{seatOnWall} build links by value (\src{model/parts/Placement.h}{127}); the default link
means ``child fore plane abuts parent aft plane'' (\src{model/parts/Placement.h}{46}). A single
resolver converts intent to geometry (\src{model/parts/Placement.cpp}{20}):

\begin{minted}{cpp}
   const double signedGap = (link.seat == SeatKind::NestInBore) ? -link.gap : link.gap;

   // Line the child's station up with the parent's, then displace by the gap. The pose locates the
   // child's fore plane, which sits c.z from its chosen station -- hence the `- c.z`.
   const double childOriginZ = p.z + signedGap - c.z;
\end{minted}

\noindent Both stations are read at live lengths and composition is coaxial in 3-DOF
(\src{model/parts/Placement.cpp}{34}); \code{Pose} pairs a fore-plane origin with a quaternion
whose \code{compose()} is the general rigid-transform law, degenerating to vector addition under
identity (\src{model/parts/Placement.h}{74}); \code{resolvePlacements} is a depth-first walk
yielding one pose per part in deterministic attachment order
(\src{model/parts/Placement.cpp}{38}).

\begin{designrule}[Placement is derived, never stored]
Absolute pose exists only as resolver output --- ``the single authority for absolute placement:
intent (StationLink) in, geometry (Pose) out. Both the simulator and the visualizer consume its
output, so the two cannot disagree'' (\src{model/parts/Placement.h}{148}). The serializer honors
the rule: a \code{.qrd} file stores only \code{<link>} intent and re-resolves on load
(\src{model/DesignSerializer.cpp}{93}; \cref{sec:persistence}).
\end{designrule}

Two caveats temper the picture. The stations are documented against \emph{live} length ---
``editing a part's length moves the seam'' (\src{model/parts/Placement.h}{42}), echoed by the
placement-gate comment's ``or a length changed'' (\src{model/parts/Part.cpp}{162}) --- yet no
geometry mutation path exists, so the live-length machinery is future-proofing with no current
trigger. And the authoring surface is far narrower than the model: the REPL's \code{addpart}
attaches everything with \cppinline|abut(offset.z())| (\src{cli/Repl.cpp}{872}), leaving
\code{NestInBore}, \code{OnSurface}, and non-default stations reachable only from a design file or
C++ (\src{model/DesignSerializer.cpp}{186}; \cref{sec:interfaces}).

\subsection{Envelope-sweep diagnostics and the fail-closed gate}\label{sec:model-sweep}

Of the two diagnostic layers shipped, exactly one is wired in. Layer~2, \code{sweepOverlaps}, is
the enforced gate: it builds a world axial interval per part from \code{axialLength()}
(\src{model/parts/Placement.cpp}{112}), excludes an offender's legitimate neighbours --- itself,
its descendants, its direct seat parent (\src{model/parts/Placement.cpp}{124}) --- sorts intervals
by aft station (\src{model/parts/Placement.cpp}{151}), samples each offender's outer radius at
feature breakpoints (its own endpoints plus every other endpoint within its span,
\src{model/parts/Placement.cpp}{160}), tests it against the smallest-id covering host's capacity
under the solid-host rule (\src{model/parts/Placement.cpp}{185}), and keeps the worst violation
per (offender, host) pair (\src{model/parts/Placement.cpp}{201}). A hollow-host guard suppresses
the false positive of an offender sitting on or outside the host skin, e.g.\ a fin set's body disc
over a co-radial aft coupler (\src{model/parts/Placement.cpp}{193}); diagnostics name parts by
stable id, never by name or pointer (\src{model/parts/Placement.h}{99}). Resolve and verdict are
cached per structural change --- mass edits never set \code{placementDirty}
(\src{model/parts/Part.cpp}{160}) --- and consumed fail-closed: \code{computeCompositeAt} refuses
a failed solve with a \cppinline|std::runtime_error| carrying the first located diagnostic, on the
stated principle that ``a self-intersecting design is a hard, located failure, never a
silently-wrong mass/inertia'' (\src{model/parts/Part.cpp}{182}), and the visualizer reads the same
cached verdict to flag the offender (\src{visualizer/RocketMesh.cpp}{440}). The sweep runs once
per edit, so the per-step cost is a flag read (\src{model/tests/DiagnosticsGateTests.cpp}{66}).

Three things bound what the sweep checks. The first is a correctness hole:

\begin{hardfail}[The monotonicity premise fails for the sphere]
The breakpoint scheme is justified in-source by ``profiles are closed-form and monotone between
breakpoints, so this is exact, not sampled'' (\src{model/parts/Placement.cpp}{161}). That premise
is false for \code{HollowSphere}, whose silhouette $\sqrt{r_o^2 - (z + r_o)^2}$ is zero at both
poles --- its own interval endpoints --- and maximal at the mid-span equator
(\src{model/parts/HollowSphere.cpp}{53}). A sphere nested inside a tube narrower than its diameter
evaluates radius $\approx 0$ at every breakpoint (unless another part's endpoint happens to land
near the equator) and passes the sweep undetected. No test exercises a sphere in the sweep: the
fixtures are tubes, cones, and fins (\src{model/tests/ResolverSweepTests.cpp}{188}). \fail
\end{hardfail}

The second is a scoping rule: at each breakpoint the sweep tests capacity against a \emph{single}
host --- the smallest id among covering, non-excluded intervals
(\src{model/parts/Placement.cpp}{185}) --- so a violation against a larger-id covering host goes
unevaluated at that station. The tie-break is pinned by test as intended
(\src{model/tests/ResolverSweepTests.cpp}{325}), but it makes the sweep a per-breakpoint
single-host check, not an all-hosts check. The third is the dormant Layer~1:
\code{radialSeamCheck} validates one seam's mated radii per \code{SeatKind}
(\src{model/parts/Placement.cpp}{56}) and is implemented and unit-tested, but nothing in
production calls it --- not \code{ensurePlacementCache}, not \code{addChildPart}, not the
serializer or CLI (\src{model/parts/Part.cpp}{170}; call sites only in
\src{model/tests/ResolverSweepTests.cpp}{194}). Its \code{zWorld} is even left parent-local behind
a comment promising ``the root-frame z is added once a pose is known''
(\src{model/parts/Placement.cpp}{87}) --- a promise no caller exists to fulfil. This is the
built-but-unconsumed motif in miniature.

\subsection{The two-pass composite walk and its caches}\label{sec:model-composite}

The tensor convention drives everything else: the bare tensor (\code{getI}/\code{setI}) is
\emph{per-unit-mass} (geometric, $\mathrm{m^2}$) for this part alone, while \code{getCompositeI}
is the full mass-weighted tensor ($\mathrm{kg\,m^2}$) of the sub-tree about the composite CM
(\src{model/parts/Part.h}{30}); the constructor seeds the composite as \cppinline|m * I|
(\src{model/parts/Part.cpp}{41}). All closed-form tensors in \code{model::InertiaTensors} ---
spheres, tube, solid cone, conical shell, and the trapezoidal fin set with its baked-in $N$-fin
rotate-and-sum (\srcf{model/InertiaTensors.h}, \src{model/InertiaTensors.cpp}{11}) --- follow the
per-unit-mass convention and are pinned against numeric-integration oracles
(\src{model/tests/InertiaTensorsTests.cpp}{186}; \cref{sec:testing}).

\code{computeCompositeAt(t)} re-weights the cached geometry by \code{getMass(t)}. Pass~1 forms the
mass-weighted composite CM, locating each part at its resolved fore-plane origin plus the local CM
station $-L/2 + \code{getCenterMassOffset()}.z()$ (\src{model/parts/Part.cpp}{199}), with a guard
that leaves a fully massless sub-tree's CM at the origin rather than dividing by zero
(\src{model/parts/Part.cpp}{210}). Pass~2 shifts each part's mass-weighted tensor to that CM with
the parallel-axis displacement tensor $f(d) = (d \cdot d)\,I_3 - d\,d^{\mathsf T}$
(\src{model/parts/Part.cpp}{19}), accumulating
\cppinline|c.mass * c.part->getI() + c.mass * parallelAxisTerm(d)|
(\src{model/parts/Part.cpp}{221}). What pass~2 deliberately omits is the $R\,I\,R^{\mathsf T}$
rotation of each child tensor --- identity in 3-DOF, skipped ``to stay bit-stable'' with an
explicit \code{TODO(6-DOF)} (\src{model/parts/Part.cpp}{217}); the class doc names the inertia
shift as ``the one piece still assuming identity'' (\src{model/parts/Part.h}{38}). The omission
already shapes the catalog: an azimuthally-arrayed component cannot be $N$ rotated children, which
is precisely why \code{FinSet} bakes its $N$-fin sum into one analytic leaf
(\src{model/parts/FinSet.h}{25}).

The result sits behind two gates: the structural flag (set by \code{setMass}/\code{setI} and the
tree mutators, propagated up the parent chain, \src{model/parts/Part.h}{239}), and the
\emph{mass-delta gate} (\src{model/parts/Part.cpp}{136}):

\begin{minted}{cpp}
   const double mNow = getCompositeMass(t);
   if(needsRecomputing || mNow != builtAtCompositeMass)
   {
      const CompositeProperties c = computeCompositeAt(t);
      compositeMass          = c.mass;
      compositeCm            = c.cm;
      compositeInertiaTensor = c.inertia;
      builtAtCompositeMass   = mNow;
      needsRecomputing       = false;
   }
\end{minted}

\noindent The exact \cppinline|!=| against a NaN-seeded sentinel (\src{model/parts/Part.h}{248})
is deliberate: during a burn the motor mass differs every step, so CM and tensor recompute every
step; after burnout \code{getMass(t)} repeats bit-for-bit and the cache freezes, immune to the
integrator's repeated and rejected stage-time queries. \code{getCompositeMass(t)} itself stays an
ungated cheap live sum --- the ODE divisor and the gate key, doing no tensor or placement work
(\src{model/parts/Part.cpp}{118}) and deliberately not diagnostics-gated
(\src{model/tests/DiagnosticsGateTests.cpp}{81}).

\begin{keyinsight}[The gate keys on total mass alone]
The mass-delta gate compares only the composite total (\src{model/parts/Part.cpp}{137}). A second
time-varying-mass part whose burn kept the total constant while shifting the CM would be served a
stale CM and tensor. Under the single-motor assumption (\cref{sec:model-rocketmodel}) this cannot
arise, but the gate's soundness is coupled to that assumption, and no test covers two simultaneous
burners.
\end{keyinsight}

One adjacent capability merits a single paragraph here. \code{getCompositeAero(refArea)} folds
each part's Barrowman contribution additively over the resolved placement, re-expressing every
$x_{cp}$ onto the sub-tree-root (tip) datum --- the same datum as \code{getCompositeCm}, so
$cp - cg$ is directly the static margin (\src{model/parts/Part.cpp}{231}) --- with the reference
area taken as the widest frontal disc, not a sum and not inflated by fins
(\src{model/parts/Part.h}{166}). The machinery is implemented and well tested at the model layer,
but its only callers are tests; the flight path does not consume it (\cref{sec:aero}).

\subsection{The concrete part catalog}\label{sec:model-catalog}

Five concrete types exist (\src{model/parts/Parts.h}{4}); \cref{tab:model-catalog} summarizes
them. Construction is mediated by \code{PartFactory}: one type-agnostic optional-field struct
\code{PartParams} is shared by the CLI parser, \code{makePart()}, and \code{params()} reflection
so the field vocabulary cannot drift (\src{model/parts/PartFactory.h}{27}); \code{makePart}
dispatches on the exact \code{typeName()} string and throws on missing fields or unknown types
(\src{model/parts/PartFactory.cpp}{27}).

\begin{table}[htbp]
\centering\small
\begin{tabularx}{\linewidth}{@{}l X X@{}}
\toprule
Part & Geometry and mass model & Distinctives and limits \\
\midrule
\code{BodyTube} &
Uniform-density hollow cylinder, $m = \rho\,\pi(r_o^2 - r_i^2)L$
(\src{model/parts/BodyTube.cpp}{28}); solid rod when $r_i = 0$
(\src{model/parts/BodyTube.h}{55}). &
Ctor validates $0 \le r_i < r_o$, $L > 0$, $\rho > 0$ (\src{model/parts/BodyTube.cpp}{22});
drag contribution left 0, skin friction named ``its eventual contribution''
(\src{model/parts/BodyTube.cpp}{38}). \\
\addlinespace
\code{ConicalNoseCone} &
Solid cone or thin lateral shell (\src{model/parts/ConicalNoseCone.h}{21}); CM at $L/4$ (solid) or
$L/3$ (shell) from the base (\src{model/parts/ConicalNoseCone.cpp}{49}). &
``A straight right circular cone'' is the only profile (\src{model/parts/ConicalNoseCone.h}{22})
--- no ogive, parabolic, elliptical, power-series, or Haack shapes. \\
\addlinespace
\code{FinSet} &
$N \ge 3$ identical trapezoidal flat-plate fins as \emph{one} analytic leaf
(\src{model/parts/FinSet.h}{22}); axial extent equals root chord
(\src{model/parts/FinSet.h}{61}). &
$N < 3$ warns and reuses the $N \ge 3$ approximation rather than throwing
(\src{model/parts/FinSet.cpp}{36}); reference area is the body disc $\pi r_b^2$, not the fin
extent (\src{model/parts/FinSet.h}{71}); only trapezoidal planforms. \\
\addlinespace
\code{HollowSphere} &
Thick-walled shell centered at $z = -r_o$, spanning $[-2r_o, 0]$
(\src{model/parts/HollowSphere.h}{51}). &
Overrides \code{isSolid()} --- the bore vanishes at the poles
(\src{model/parts/HollowSphere.h}{59}); doubles as the boot placeholder body
(\src{model/RocketModel.cpp}{27}). \\
\addlinespace
\code{Motor} &
Zero-length leaf wrapping a \code{MotorModel} by value; \code{getMass(t)} returns the burn-time
mass (\src{model/parts/Motor.h}{40}); tensor is a constant solid cylinder, or zero when geometry
is unknown (\src{model/parts/Motor.cpp}{25}). &
The one override that makes the tree's mass, CG, and inertia time-aware
(\src{model/parts/Motor.h}{18}). No geometric envelope, so motor fit is never checked by the
sweep (\src{model/RocketModel.cpp}{118}); not factory-constructible
(\src{model/parts/PartFactory.cpp}{73}). \\
\bottomrule
\end{tabularx}
\caption{The concrete part catalog at commit \code{254cd41}.}\label{tab:model-catalog}
\end{table}

Measured against what a real rocket needs, the catalog is sparse. Missing outright are transitions
and boattails, curved nose profiles, inner and motor-mount tubes, tube couplers as a first-class
type (a coupler today is a nested \code{BodyTube}), centering rings, bulkheads, launch lugs and
rail buttons, recovery devices (parachutes, streamers, shock cords), discrete mass objects, and
pods or strap-on boosters (\src{model/parts/PartFactory.cpp}{75}). \fabsent{} One convention is
worth flagging as convention rather than constraint: every committed \code{.qrd} fixture mounts
its \code{FinSet} with \code{seat="OnSurface"} on a \code{BodyTube} (e.g.\
\src{tests/data/designs/large38\_basic.qrd}{13}), but nothing enforces this --- a REPL-authored
fin set gets an \code{Abut} link. The quantitative OpenRocket comparison is consolidated in
\cref{sec:comparison}; the ranked gap list, in \cref{sec:roadmap}.

\subsection{RocketModel: tree root and the Propagatable bridge}\label{sec:model-rocketmodel}

\code{RocketModel} holds the tree root and implements \code{Propagatable}, the
model$\leftrightarrow$sim bridge (\src{model/Propagatable.h}{23}, \src{model/RocketModel.h}{31}).
On the mass side --- this section's half of the bridge --- it is a thin, honest delegation:
\code{getMass(t)} returns \cppinline|topPart->getCompositeMass(t)| with the motor included as a
child part, \code{setMass} writes only the top part's dry mass (\src{model/RocketModel.cpp}{36},
\src{model/RocketModel.h}{87}), \code{getCompositeInertiaTensor} delegates to
\code{getCompositeI} (\src{model/RocketModel.cpp}{43}), and \code{writeMassProperties} snapshots
the mutually consistent gated mass/CG/inertia triple into \code{StateData} per recorded step
(\src{model/RocketModel.cpp}{53}; called at \src{sim/Propagator.cpp}{132}).
\code{terminateCondition} ends the flight when descending below the launch site, $z < 0$ and
$v_z < 0$ (\src{model/RocketModel.cpp}{62}). The force side --- world-$z$ thrust ``assumed through
the CM'' (\src{model/RocketModel.cpp}{70}), a flat $C_d A$ drag term, \code{getTorques} returning
zero (\src{model/RocketModel.cpp}{94}) --- belongs to \cref{sec:aero}; the standing consequence
here is that the composite $I(t)$ and $CG(t)$ the tree so carefully maintains have, today, exactly
one consumer: the per-step state snapshot.

The part-tree facade keeps a borrowed \cppinline|part::Motor*| handle consistent across edits:
\code{setRoot}, \code{addPart}, and \code{removePart} all re-resolve it by a depth-first search
for the \emph{first} \code{Motor} node --- an explicit single-motor assumption
(\src{model/RocketModel.cpp}{13}); a second motor in the tree would be silently ignored by thrust.
\code{setMotorModel} creates the \code{Motor} child on first call and swaps the wrapped model
in place on later calls, so the borrowed pointer never dangles (\src{model/parts/Motor.cpp}{16});
where that child seats --- at the airframe CM with a hard-coded zero offset, one motor, no stages
--- is \cref{sec:motors}'s subject. Two smaller dormant seams
round out the picture: \code{deriveReferenceAreaFromGeometry()} computes the Barrowman-convention
reference area from the tree (\src{model/RocketModel.cpp}{48}) but has no production caller ---
the drag model runs on a hard-coded 38\,mm-tube default (\src{model/RocketModel.h}{141};
\cref{sec:aero}) --- and \code{Propagatable} still carries a default-constructed
\code{sim::Aero aeroData} member that nothing reads (\src{model/Propagatable.h}{61},
\src{sim/Aero.h}{63}). Both fit the layer's character: the seams toward 6-DOF and real
aerodynamics are built and tested here first, and \cref{sec:roadmap} ranks what it will take to
consume them.
